(a) Put $R=k[t]$. Every submodule of $R^m$ is free: project to the first coordinate, whose image is a principal ideal. A lift of its generator splits off a free summand, and induction applies to the kernel inside $R^{m-1}$. Thus every finite $R$-module $M$ has a resolution $0\to R^s\to R^r\to M\to0$, giving $[M]=(r-s)[R]=\operatorname{rank}(M)[R]$. Under the correspondence between finite modules and coherent sheaves, rank therefore identifies $K(\mathbb A^1_k)$ with $\mathbb Z$.

(b) Taking the generic stalk preserves exact sequences, and dimensions of finite-dimensional $K(X)$-vector spaces are additive. Thus rank defines a homomorphism $K(X)\to\mathbb Z$, which is surjective since $\operatorname{rank}\mathcal O_X=1$.

(c) Let $i:Y\to X$ and $j:U=X-Y\to X$ be the inclusions. The functors $i_*$ and $j^*$ are exact and preserve coherence, and $j^*i_*=0$. Let $H\subseteq K(X)$ be generated by the classes of coherent sheaves supported on $Y$. If $\mathcal I$ defines $Y$ and $\operatorname{Supp}\mathcal F\subseteq Y$, coherence and quasi-compactness give $\mathcal I^N\mathcal F=0$ for some $N$. The filtration by $\mathcal I^m\mathcal F$ has coherent $\mathcal O_Y$-module quotients. Hence $H=\operatorname{im}K(Y)$.

By \exref{II.5.15}, every coherent sheaf on $U$ extends to one on $X$. Any two extensions $\mathcal F,\mathcal G$ of the same sheaf have the same class modulo $H$: extend the graph of their identification on $U$ to a coherent subsheaf $\mathcal C\subseteq\mathcal F\oplus\mathcal G$ using \exref{II.5.15}(d). The kernels and cokernels of both projections from $\mathcal C$ are supported on $Y$, so $[\mathcal F]=[\mathcal C]=[\mathcal G]$ modulo $H$.

For an exact sequence on $U$, extend its middle term, extend its first term as a coherent subsheaf of that extension, and take the quotient. This gives an exact sequence on $X$ restricting to the given one. Consequently the class modulo $H$ of an extension defines an inverse $K(U)\to K(X)/H$ to restriction. This proves the required exact sequence.
